\setcounter{equation}{0}
\section{Simulation}

This project aims to produce a quantum-computer simulator that runs on classical computers. Unfortunately, any classical simulator is fundamentally inefficient compared to the quantum hardware it models. Although a quantum computer operates on a different physical paradigm than a classical computer, quantum computers cannot compute anything that classical computers cannot --- the gain is in efficiency, not in the set of computable functions. Classical machines can simulate quantum systems to arbitrary precision by storing the state-vector amplitudes and applying the unitary transformations the algorithm prescribes.

Using a classical system to simulate a quantum computer nullifies any speed advantage of the quantum algorithm: most of the structure that quantum computers exploit implicitly must be programmed and manipulated explicitly. Simulating even a moderate quantum algorithm of, say, $100$ qubits would require storing $2^{100}$ (roughly $10^{30}$) complex amplitudes --- more than any conceivable classical machine can hold, and even more far beyond what one could rotate by an arbitrary unitary in any reasonable time.

So although a quantum computer can in principle be simulated by a classical computer, the simulation becomes exponentially expensive as the qubit count grows. The memory required to store the state vector and (depending on the formulation) the operator matrices grows exponentially with the number of qubits, which is the bottleneck the parallelisation strategy of the next section is designed to push back as far as possible.
